% Manuscript skeleton for the HiGS training-speed submission.
%
% Status: DRAFT - no confirmatory claims are asserted yet. Every number must
% trace to paper/claims.yaml after the frozen confirmatory matrix completes.
% Replace this preamble with the target venue template before submission.
\documentclass[11pt,letterpaper]{article}

\usepackage[utf8]{inputenc}
\usepackage[T1]{fontenc}
\usepackage{microtype}
\usepackage{booktabs}
\usepackage{graphicx}
\usepackage{amsmath}
\usepackage{amssymb}
\usepackage{hyperref}
\usepackage{xcolor}

\title{Trainable Hierarchical Gaussian Splatting:\\Skipping Irrelevant Work During Training}
\author{%
  \textbf{TODO: author list and affiliations}\\*
  \texttt{TODO: contact email}
}
\date{\today}

\newcommand{\TODO}[1]{\textcolor{red}{\textbf{[TODO] }#1}}

\begin{document}

\maketitle

\begin{abstract}
\TODO{Write after the confirmatory matrix (30k steps, held-out scenes, paired
seeds, block-bootstrap intervals) is complete. Current draft claims are
exploratory and must not be presented as final results.}
\end{abstract}

\section{Introduction}

\label{sec:intro}

\TODO{State the problem: end-to-end 3D Gaussian Splatting training cost.
Position the contribution as a trainable hierarchical representation that
spends compute only on currently relevant geometry. Explicitly separate
upstream components (gsplat, Official 3DGS), this repository's
implementation contribution, the frozen method, exploratory negative results,
and confirmatory results.}

\section{Related Work}

\label{sec:related}

\TODO{Cover differentiable Gaussian rasterization, hierarchical and sparse
training, tile-sampled training, progressive-resolution schedules, and
per-Gaussian culling. Keep primary-source citations in
\texttt{references.bib} and only cite claims verified against the pinned
artifacts.}

\section{Method}

\label{sec:method}

\subsection{Overview}

\TODO{One figure showing the training pipeline: progressive-resolution
schedule, error-guided tile sampling, projection-only visibility refresh,
masked-Adam updates, and cull-mask caching.}

\subsection{Error-guided sampling}

\TODO{Describe the sampling rule and why tile granularity creates correlation
noise at high Gaussian counts (rounds 51, 53, 54).}

\subsection{Visibility refresh and masked updates}

\TODO{Describe the camera-set-keyed cull-mask cache (round 59) and the
masked-Adam kernel (round 60). Report correctness verification: masked rows
bit-identical, unmasked rows within 1--2 float32 ULPs.}

\subsection{Frozen configuration}

\TODO{List the exact flags and hyperparameters of the frozen speed and
quality configurations once confirmatory evaluation starts. Until then this
subsection stays empty: exploration is not a method.}

\section{Experiments}

\label{sec:experiments}

\subsection{Benchmark protocol}

All Tier~A measurements use protocol
\texttt{matrix-primary-1080p-v2}
(SHA-256 \texttt{892e18890501c408dc6746af69f17e16973604f0c02a3caddf1954d3bf1fede2})
on a single NVIDIA A100-SXM4-80GB host (EPIC-05) at 1920$\times$1080.
Each case runs 100 evaluated frames per repeat and five repeats. Quality
metrics are PSNR, SSIM, and LPIPS against the canonical ground-truth crop.

\subsection{Datasets}

Five canonical scenes: Mip-NeRF 360 (garden, bicycle, bonsai) and
Tanks and Temples (truck, train). \TODO{Add a held-out dataset family not used
in exploratory rounds before any confirmatory claim.}

\subsection{Training configurations}

\TODO{Tabulate the frozen speed and quality configurations, the baseline
(upstream gsplat full-resolution), the 30k-step convergence horizon, the
paired seed set, and per-scene/per-seed wall time, time-to-target-quality,
final Gaussian count, and peak VRAM.}

\subsection{Statistical protocol}

The experimental unit is one training run (scene $\times$ seed). Paired
method-baseline deltas are summarized per scene and the scene-level bootstrap
distribution is computed with the block bootstrap in
\texttt{src/scripts/bootstrap\_analysis.py}; we report 95\% percentile
intervals and a paired effect size. \TODO{Add per-seed raw values and
pre-registered stopping rules before submission.}

\section{Results}

\label{sec:results}

\TODO{Insert tables generated from \texttt{paper/tables/}; every cell must map
to an immutable JSON artifact listed in \texttt{claims.yaml}. Current
exploratory round summaries (R56--R61) are evidence of a bounded research
space, not confirmatory results.}

\section{Limitations}

\label{sec:limitations}

\TODO{State: single GPU family; 3,000-step exploratory horizon; five scenes
used during development; host contention as a confounder in latency studies;
and any negative results that bound the method.}

\section{Reproducibility}

\label{sec:reproducibility}

The canonical suite, protocol hash, renderer pinning, and evidence tiers are
documented in \texttt{docs/reproducibility.md}. All Tier~A artifacts live
under \texttt{results/measured/} and are never mixed with diagnostic runs
under \texttt{results/diagnostics/}. \TODO{Add the clean-machine artifact
evaluation recipe and the single claim-to-JSON manifest.}

\section{Ethics and Impact}

\TODO{Brief standard statement: datasets used under their original terms (see
\texttt{paper/LICENSE-STATEMENT.md}), no human subjects, no new environmental
footprint beyond documented benchmark runs.}

\bibliographystyle{plain}
\bibliography{references}

\end{document}